\documentclass[10pt,a4paper]{amsart}
\usepackage[margin=19mm]{geometry}
\usepackage{amsmath,amssymb,mathtools}
\usepackage[scr=rsfs,cal=euler]{mathalfa}
\usepackage{xfrac}
\usepackage[unicode,hidelinks,bookmarks=false]{hyperref}
\newcommand{\ie}{i.e.\ }
\newcommand{\cf}{cf.\ }
\newcommand{\resp}{respectively\ }
\newcommand{\Subsection}{\subsection}
\providecommand{\nobreakdash}{\nobreak\hbox{-}}
\setlength{\emergencystretch}{2em}
\multlinegap0pt
\newtheorem{theorem}{Theorem}[section]
\newtheorem{lemma}[theorem]{Lemma}
\newtheorem{prop}[theorem]{Proposition}
\newtheorem{coro}[theorem]{Corollary}
\theoremstyle{definition}
\newtheorem{defi}[theorem]{Definition}
\newtheorem{rema}[theorem]{Remark}
\providecommand{\eg}{e.g.\ }
\newcounter{cdrthm}
\renewcommand{\theHcdrthm}{\thetheorem.\arabic{cdrthm}}
\newenvironment{enonce}[2][remark]{\refstepcounter{cdrthm}\par\medskip\noindent\textit{#2 \thecdrthm.}\ }{\par}
\newcommand{\un}{\underline}

\renewcommand{\epsilon}{\varepsilon}
\renewcommand{\tilde}{\widetilde}
\newcommand{\G}{\mathbb{G}}
\newcommand{\R}{\mathbb{R}}
\newcommand{\Z}{\mathbb{Z}}
\newcommand{\D}{\mathbf{D}}
\newcommand{\B}{\mathcal{B}}
\newcommand{\T}{\mathcal{T}}
\newcommand{\Sand}{\mathcal{S}}
\renewcommand{\L}{\mathcal{L}}
\renewcommand{\P}{\mathcal{P}}
\newcommand{\what}{\widehat}
\newcommand{\col}[1]{x_{#1}}
\newcommand{\hcol}[1]{\widehat x_{#1}}
\newcommand{\indB}[1]{k_{#1}}
\newcommand{\CB}{x_B}
\newcommand{\hCB}{x_{\widehat B}}
\newcommand{\colB}[1]{x_{k_{#1}}}
\newcommand{\hcolB}[1]{\widehat x_{k_{#1}}}
\newcommand{\CS}{x_S}
\newcommand{\w}{\mathfrak w}

\DeclareMathOperator{\Supp}{Supp}
\DeclareMathOperator{\Diff}{Diff}
\DeclareMathOperator{\Spec}{Spec}
\DeclareMathOperator{\Proj}{Proj}
\DeclareMathOperator{\SL}{Supp}
\DeclareMathOperator{\Sym}{Sym}
\DeclareMathOperator{\Pic}{Pic}
\DeclareMathOperator{\Ima}{im}
\DeclareMathOperator{\sgn}{sgn}
\DeclareMathOperator{\Cut}{Cocirc}
\DeclareMathOperator{\Flow}{Circ}
\DeclareMathOperator{\E}{Ext}
\DeclareMathOperator{\Ann}{Ann}
\DeclareMathOperator{\Aut}{Aut}
\DeclareMathOperator{\Br}{Br}
\DeclareMathOperator{\cd}{cd}
\DeclareMathOperator{\Char}{char}
\DeclareMathOperator{\Cl}{Cl}
\DeclareMathOperator{\codim}{codim}
\DeclareMathOperator{\Div}{Div}
\DeclareMathOperator{\rk}{rk}
\DeclareMathOperator{\Gr}{Gr}
\DeclareMathOperator{\spn}{span}

\newcommand*{\mk}{\mkern -1mu}
\newcommand*{\Mk}{\mkern -2mu}
\newcommand*{\mK}{\mkern 1mu}
\newcommand*{\MK}{\mkern 2mu}

\hypersetup{urlcolor=purple, linkcolor=blue, citecolor=red}


\newcommand*{\romanenumi}{\renewcommand*{\theenumi}{\roman{enumi}}}
\newcommand*{\Romanenumi}{\renewcommand*{\theenumi}{\Roman{enumi}}}
\newcommand*{\alphenumi}{\renewcommand*{\theenumi}{\alph{enumi}}}
\newcommand*{\Alphenumi}{\renewcommand*{\theenumi}{\Alph{enumi}}}

\title[Geometric matrix-tree multijection]{A family of matrix-tree multijections: original Theorem6.10}
\author{Alex McDonough}
\date{Source: Algebraic Combinatorics4(5)(2021),795--822; DOI10.5802/alco.181}
\begin{document}\maketitle
\noindent\textit{Editorial scope.} The counted unit is original Theorem6.10 with the author\textquotesingle s complete construction of its map. All other results are uncounted prerequisites or local core. Graph examples and later slicing/chamber interpretations are omitted. Original mathematical text and formulas are preserved. The source identifies its Theorem1.1 below as a reframing of \cite[Theorem8.1]{CutsFlows}.
\section{Original named matrix-tree prerequisite}
Let $D$ be a standard representative matrix (see Definition~\ref{def:standrep}). In Section~\ref{sec:Background}, we define the \emph{sandpile group} $\Sand(D)$, the \emph{bases} $\B(D)$, and the \emph{basis multiplicity function} $m$ which maps each $B \in \B(D)$ to a positive integer. In this context, we get the following theorem, which is a reframing of Theorem 8.1 from~\cite{CutsFlows}.

\begin{theorem}[Sandpile matrix-tree theorem on standard representative matrices]\label{thm:OAmtt} 
\[
|\Sand(D)| = \sum_{B \in \mathcal B(D)} m(B)^2.
\]
\end{theorem}

\section{Notational Conventions}\label{sec:notation}
We will write $\Z$ for the integers, $\Z_{\ge 1}$ for the positive integers, and $\R$ for the real numbers. We write $[a,b]$ for the set $\{x\in \Z\mid a \le x \le b\}$ and $[b]$ for $[1,b]$. We denote a vector of all zeros by $\un 0$. We use the symbol $D$ for an $r \times n$ integer matrix which, starting in Section~\href{https://alco.centre-mersenne.org/articles/10.5802/alco.181/}{4}, will always be a \emph{standard representative matrix} (see Definition~\ref{def:standrep}). We write $\what D$ and $\D$ for the \emph{dual matrix} and \emph{full matrix} of $D$ respectively (again, see Definition~\ref{def:standrep}). We will always write the determinant of a square matrix $A$ as $\det(A)$ and use $| \cdot |$ for set cardinality or absolute value. We also write $A^T$ for the transpose of a matrix $A$ and $I_k$ for the $k \times k$ identity matrix. We will frequently be working with polyhedra embedded in $\R^k$ (where $k$ is either $n$, $r$, or $n-r$). We use the term \emph{volume} to mean $k$-dimensional Lebesgue measure.

\section{Background and definitions}\label{sec:Background}

\begin{defi}
A \emph{lattice} is a subgroup of a finite-dimensional vector space that is isomorphic to $\Z^k$ for some $k$.
\end{defi}

A subgroup of a lattice is called a \emph{sublattice}. Given any set $S$ of $\Z^k$ vectors, the integer linear combinations of these vectors form a lattice $L$ of dimension at most $k$. We say that $S$ \emph{generates} $L$. If the vectors in $S$ are linearly independent, we say that $S$ is an \emph{integral basis} for $L$. 

\begin{rema} When working with vector spaces, any maximal linearly independent set of generators is a basis. However, a maximal linearly independent set of generators for a lattice is not always an integral basis of this lattice. For example, the set $\{2,3\}$ generates $\Z$, but neither $\{2\}$ nor $\{3\}$ is an integral basis for $\Z$. 
\end{rema}

\begin{prop}[{\cite[Theorem~14.5.3]{Godsil}}]\label{cokersize}
If $S$ is a set of $k$ vectors in $\Z^k$ that are an integral basis for a lattice $L$, then the group $\Z^k/ L$ has size equal to the magnitude of the determinant of the matrix formed by the vectors of $S$.
\end{prop}

For a lattice $L$, the group $\Z^k / L$ is called the \emph{cokernel} of $L$. For some integers $r \le n$, let $D$ be an $r \times n$ integer matrix:

\begin{defi}\label{def:cutandflow}\hfill
\begin{itemize}
\item The \emph{cocircuit space} of $D$ is the space $\Ima_\R(D^T)$.
\item The \emph{circuit space} of $D$ is the space $\ker_\R(D)$.
\item The \emph{cocircuit lattice} of $D$ is the lattice $\Ima_\Z(D^T)$.
\item The \emph{circuit lattice} of $D$ is the lattice $\ker_\Z(D)$.
\item The \emph{sandpile lattice} of $D$ is the lattice $\Ima_\Z(D^T) \oplus \ker_\Z(D)$.
\end{itemize}
\end{defi}

The cocircuit space and cocircuit lattice are generated by the rows of $D$. The circuit space and circuit lattice are generated by the coefficients of integer linear combinations of columns of $D$ that sum to $\un 0$. Note that these generators are all elements of $\Z^n$.



\setcounter{theorem}{5}
\begin{lemma}[{\cite[Proposition~5.1]{CutsFlows}}]\label{lem:orthogonality}
For any integer matrix $D$, the spaces $\Ima_\R(D^T)$ and $\ker_\R(D)$ are orthogonal complements. 
\end{lemma}

Note that because $\Ima_\Z(D^T) \subset \Ima_\R(D^T)$ and $\ker_\Z(D) \subset \ker_\R(D)$, this also means that $\Ima_\Z(D^T)$ and $\ker_\Z(D)$ are always orthogonal. We also get the following corollary:

\begin{coro}\label{cor:cutofflow}
Let $D$ be an integer matrix and $\what D$ be a matrix with rows that generate $\ker_\R(D)$. Then, $\ker_\R(D) = \Ima_\R(\what D^T)$ and $\ker_\R(\what D) = \Ima_\R(D^T)$. 
\end{coro}

\begin{proof}
The first equality follows immediately from the fact that $\Ima_\R(\what D^T)$ is generated by the rows of $\what D$ which also generate $\ker_\R(D)$ by definition. 

For the second equality, by Lemma~\ref{lem:orthogonality}, $\ker_\R(\what D)$ is the orthogonal complement of $\Ima_\R(\what D^T)$ which we established is equal to $\ker_\R(D)$. By a second application of Lemma~\ref{lem:orthogonality}, $\ker_\R(D)$ is the orthogonal complement of $\Ima_\R(D^T)$. Since the composition of two orthogonal complements is the identity, we conclude that $\ker_\R(\what D) = \Ima_\R(D^T)$. 
\end{proof}


\begin{defi}\label{def:standrep}
An $r \times n$ integer matrix $D$ is a \emph{standard representative matrix} if it is of the form:
\[
D = \begin{pmatrix} I_r & M\end{pmatrix},
\]
where $I_r$ is the $r \times r$ identity matrix and $M$ is any $r \times (n-r)$ integer matrix. A standard representative matroid $D$ is associated with two other matrices:
\[
\what D = \begin{pmatrix} -M^T & I_{n-r}\end{pmatrix} \hspace{.4 cm}\text{ and }\hspace{.4 cm} \D = \begin{pmatrix} D \\ \what D \end{pmatrix} = \begin{pmatrix} I_r & M\\ -M^T & I_{n-r}\end{pmatrix}.
\]
We call $\what D$ the \emph{dual matrix} of $D$ and $\D$ the \emph{full matrix} of $D$. We will show in Lemma~\ref{lem:flowofcut} that our notation for $\what D$ is consistent with Corollary~\ref{cor:cutofflow}. 
\end{defi}

\setcounter{theorem}{9}
\begin{lemma}[{\cite[Corollary~4.6.6]{mythesis}}]\label{lem:flowofcut}
If $D$ is a standard representative matrix, then $\ker_\Z(D) = \Ima_\Z(\what D^T)$, $\Ima_\Z(D^T) = \ker_\Z(\what D)$, and $\Ima_\Z(\D^T) = \Ima_\Z(D^T) \oplus \ker_\Z(D)$. 
\end{lemma}

\begin{defi}\label{def:sandpile}
The \emph{sandpile group} of a standard representative matrix $D$, denoted $\Sand(D)$, is the finite abelian group
\[
\Sand(D) = \Z^n / (\Ima_\Z(D^T) \oplus \ker_\Z(D)) = \Z^n / (\Ima_\Z(\D^T)).
\]
\end{defi}

Notice that by Lemma~\ref{lem:flowofcut}, the sandpile group of $D$ is the cokernel of the sandpile lattice of $D$. 

\begin{defi}\label{def:bandm}
The set of \emph{bases} of $D$, written $\mathcal B(D)$, is the set of $r$-tuples of columns of $D$ such that the determinant of $D$ restricted to these columns is nonzero. For $B \in \mathcal B(D)$, let $m(B)$ be the absolute value of this determinant. This $m(B)$ is called the \emph{multiplicity} of $B$. 
\end{defi}

\setcounter{theorem}{14}
\begin{defi}
Let $S$ and $T$ be sets and $\mathfrak m$ be a map from $T \to \Z_{\ge 1}$. An \emph{$\mathfrak m$-multijection} between $S$ and $T$ is a map $f:S \to T$ such that for all $t \in T$, $|f^{-1}(t)| = \mathfrak m(t)$. \end{defi}

An $\mathfrak m$-multijection can also be thought of as a bijection between $S$ and a multiset consisting of $\mathfrak m(t)$ copies of each $t \in T$. In this paper, we give an explicit procedure for constructing $\mathfrak m$-multijections between $\Sand(D)$ and $\B(D)$ for $\mathfrak m(B) = m(B)^2$. To do this, we use a geometric construction, which also produces a periodic tiling of $\R^n$. 

\setcounter{section}{4}
\section{A Tiling of \texorpdfstring{$\R^n$}{Rn}}\label{sec:tiling}

For the remainder of this paper, we will always let $D = \begin{pmatrix} I_r &M \end{pmatrix}$ be an $r \times n$ standard representative matrix. Furthermore, we let $\widehat D = \begin{pmatrix} -M^T & I_{n-r} \end{pmatrix}$ be the dual matrix of $D$ and 
\[
\D = \begin{pmatrix} D \\ \what D \end{pmatrix} = \begin{pmatrix} I_r & M\\ -M^T & I_{n-r}\end{pmatrix}
\]
be the full matrix of $D$. Recall from Definition~\ref{def:bandm} that $\B(D)$ is the set of $r$ element subsets of the columns of $D$ with nonzero determinant and for $B \in \B(D)$, $m(B)$ is the magnitude of the corresponding determinant. In this section, we will associate each $B \in \B(D)$ with a lattice parallelepiped and then show that the non-convex polytope formed by their union periodically tiles $\R^n$. In the next section, we will show how to use this tiling to construct a family of multijections. 

We think of $B \in \B(D)$ as a set $\{\indB 1,\dots, \indB r\}$ of column indices. These simultaneously describe a set of columns of $D$, $\widehat D$ or $\D$. Because we are working in $\R^n$, it will be useful to allow for a version of the sandpile group whose representatives are real vectors. 

\begin{defi}
The \emph{continuous sandpile group} of $D$ is the group:
\[
\tilde \Sand(D) = \R^n / (\Ima_\Z(D^T) \oplus \ker_\Z(D)) = \R^n/\Ima_\Z(\D^T).
\]
\end{defi}

We will also make heavy use of the following lemma, which follows immediately from the definition of sandpile groups and continuous sandpile groups of standard representative matrices. 
\begin{lemma}\label{lem:whenequiv}
Let $D$ be an $r \times n$ standard representative matrix. Two vectors $z,z' \in \Z^{n}$ (resp. $\R^{n}$) are equivalent as elements of $\Sand(D)$ (resp. $\tilde \Sand(D)$) if and only if $z - z' \in \Ima_\Z(\D^T)$.
\end{lemma}

We introduce some definitions and notation that can be found in~\cite{Beck}.
\begin{defi}\hfill\label{def:parpip}
\begin{itemize}
 \item The \emph{fundamental parallelepiped} of a square matrix $A$ with column vectors $\{\col 1,\dots,\col k\}$ is the set of points:
\[
 \left\{\sum_{i=1}^k a_i \col i \mid 0 \le a_i \le 1\right\}.
\]
 \item The \emph{half-open fundamental parallelepiped} of a square matrix $A$ with column vectors $\{\col 1,\dots,\col k\}$ is the set of points:
\[
 \left\{\sum_{i=1}^k a_i \col i \mid 0 \le a_i < 1\right\}.
\]
\end{itemize}
We use the notation $\Pi_{\bullet}(A)$ to indicate the fundamental parallelepiped of $A$ and 
$\Pi_{\circ}(A)$ to indicate the half-open fundamental parallelepiped of $A$. 
\end{defi}

It is a classical result that the volume of $\Pi_\bullet(A)$ or $\Pi_\circ(A)$ is the magnitude of $\det(A)$. 

\begin{defi} \label{def:p1p2p}
For any basis $B\in \B(D)$:
\begin{itemize}
 \item $P_1(B)$ is the fundamental parallelepiped of $D$ restricted to columns in $B$. 
 \item $P_2(B)$ is the fundamental parallelepiped of $\widehat D$ restricted to columns \emph{not} in $B$.
 \item $P(B)$ is the direct product of $P_1(B)$ and $P_2(B)$.
\end{itemize}
\end{defi}

Note that $P_1(B)$ is $r$-dimensional, $P_2(B)$ is $(n-r)$-dimensional, and $P(B)$ is $n$-dimensional. 

\begin{lemma}[{\cite[Lemma~7.1.5]{mythesis}}]\label{lem:pipedvol}For any basis $B\in \mathcal B(D)$, $P_1(B)$ and $P_2(B)$ each have volume $m(B)$ while $P(B)$ has volume $m(B)^2$.
\end{lemma}

We can also describe $P(B)$ in the following way.

For each column of $\D$, if this column corresponds to an index of $B$, replace the last $(n-r)$ entries with 0's. If this column does not correspond to an index of $B$, replace the first $r$ entries with 0's. The fundamental parallelepiped of this matrix is $P(B)$. 

\setcounter{theorem}{6}
\begin{prop}\label{prop:nointersect}
The parallelepipeds $P(B)$ for each basis $B\in \B(D)$ do not intersect except at their boundaries.
\end{prop}
\begin{proof}
Let $\{\col 1,\dots, \col n\}$ be the columns of $D$ and $\{\hcol 1,\dots, \hcol n\}$ be the columns of $\widehat D$. Let $B_1$ and $B_2$ be two distinct bases in $\B(D)$. 

$P(B_1)$ and $P(B_2)$ have intersecting interiors if and only if $P_1(B_1)$ and $P_1(B_2)$ have intersecting interiors and $P_2(B_1)$ and $P_2(B_2)$ have intersecting interiors. Assume that $P_1(B_1)$ and $P_1(B_2)$ have intersecting interiors. Then, for some coefficients $a_1,\dots,a_{n},\\
b_1,\dots,b_{n}$, we have the following equality
\[
\sum_{i=1}^{n} a_i \col i = \sum_{i=1}^{n} b_i \col i
\]
where $0 < a_i < 1$ for $i \in B_1$, $a_i = 0$ for $i \not\in B_1$, $0 < b_i < 1$ for $i \in B_2$, and $b_i =0 $ for $i \not\in B_2$. If we subtract the second sum from the first and define $d_i = a_i - b_i$, we get the equation
\[
\sum_{i=1}^{n} d_i \col i=\un 0.
\]

The above equation implies that
\[
(d_1,d_2,\dots,d_n)\in \ker_\mathbb R (D).
\]

Similarly, if $P_2(B_1)$ and $P_2(B_2)$ have intersecting interiors then for some coefficients $\widehat a_1,\dots,\widehat a_{n},
\widehat b_1,\dots,\widehat b_{n}$, we have the following equality:
\[
\sum_{i=1}^{n} \widehat a_i\hcol i = \sum_{i=1}^{n} \widehat b_i\hcol i
\]
where $0 < \widehat a_i <1$ for $i \not\in B_1$, $a_i = 0$ for $i \in B_1$, $0 < b_i < 1$ for $i \not\in B_2$, and $b_i =0 $ for $i \in B_2$. For each $i \in [n]$, let $\widehat d_i = \widehat a_i - \widehat b_i$. Then,
\[
\sum_{i=1}^{n} \widehat d_i\hcol i=\un 0.
\]

It follows that:
\[
(\widehat d_1,\widehat d_2,\dots,\widehat d_n)\in \ker_\mathbb R (\widehat D) = \Ima_\R(D^T)
\]
where the last equality follows from Lemma~\ref{lem:flowofcut}.

Lemma~\ref{lem:orthogonality} says that $\Ima_\R(D^T)$ and $\ker_\R(D)$ are orthogonal. This means,
\[
0 = (d_1,d_2,\dots,d_n) \cdot(\widehat d_1,\widehat d_2,\dots,\widehat d_n) = \sum_{i=1}^{n} d_i\widehat d_i.
\]

For each $i$, there are 4 possibilities:

\begingroup
\setcounter{cdrthm}{0}
\renewcommand\thecdrthm{\arabic{cdrthm}}

%%%1
\begin{enonce}[remark]{Case}\label{proof_def5.7_cas1} $i \in B_1 \cap B_2$:

$\widehat a_i = \widehat b_i = 0$, so $\widehat d_i = 0$ and $d_i \cdot \widehat d_i = 0$. 
\end{enonce}

%%%2
\begin{enonce}[remark]{Case}\label{proof_def5.7_cas2} $i \in B_1 \setminus B_2$:

$b_i = 0$, so $d_i = a_i$. This means that $0 < d_i < 1$. Furthermore, $\widehat a_i = 0$, so $\widehat d_i = -\widehat b_i$ and $-1 < \widehat d_i < 0$. It follows that $d_i \cdot \widehat d_i < 0$.
\end{enonce}

%%%3
\begin{enonce}[remark]{Case}\label{proof_def5.7_cas3} $i \in B_2 \setminus B_1$:

$a_i = 0$, so $d_i = -b_i$. This means that $-1 < d_i < 0$. Furthermore, $\widehat b_i = 0$, so $\widehat d_i = \widehat a_i$ and $0 < \widehat d_i < 1$. It follows that $d_i \cdot \widehat d_i < 0$.
\end{enonce}

%%%4
\begin{enonce}[remark]{Case}\label{proof_def5.7_cas4} $i \not\in B_1 \cup B_2$:

$ a_i = b_i = 0$, so $d_i = 0$ and $d_i \cdot \widehat d_i = 0$.\\

$B_1$ and $B_2$ are the same size and distinct, so cases~\ref{proof_def5.7_cas2} and~\ref{proof_def5.7_cas3} must each occur at least once. This means that
\[
\sum_{i=1}^{n} d_i\widehat d_i < 0.
\]
This is a contradiction.\qedhere
\end{enonce}
\endgroup
\setcounter{cdrthm}{7}
\end{proof}

\begin{defi}
$T(D)$, the \emph{tile associated with $D$}, is 
\[
T(D) = \bigcup_{B \in \B(D)} P(B).
\]
\end{defi}

%\noindent 
Corollary~\ref{cor:tiling} will justify why we call this non-convex polyhedron a \emph{tile}. 

The following corollary follows directly from Lemma~\ref{lem:pipedvol} which gives the size of each $P(B)$ and Proposition~\ref{prop:nointersect} which says that they don't intersect.

\begin{coro}\label{cor:voltile}
The volume of $T(D)$ is equal to
\[
\sum_{B \in \B(D)} m(B)^2.
\]
\end{coro}

%\noindent 
Note that this sum is also equal to $|\Sand(D)|$ by Theorem~\ref{thm:OAmtt}. 

When considering all of $T(D)$, we can strengthen Proposition~\ref{prop:nointersect} to the following:

\begin{prop}\label{prop:lastDomino}
Two distinct points of $T(D)$ can only be equivalent as elements of $\tilde \Sand(D)$ if they are both on the boundary of $T(D)$. 
\end{prop}

\begin{proof}
First, we show that two points of $T(D)$ can only be equivalent as elements of $\tilde \Sand(D)$ if they are each on the boundary of some $P(B)$. 

For some $B_1,B_2 \in \B(D)$, let $p_1$ and $p_2$ be interior points of $P(B_1)$ and $P(B_2)$ respectively. Using the notation and reasoning from Proposition~\ref{prop:nointersect}, we can write $p_1 - p_2$ as the vector whose first $r$ entries are given by 
\[
\sum_{i=1}^{n} d_i\col i,
\]
and whose last $n-r$ entries are given by 
\[
\sum_{i=1}^{n} \widehat d_i\hcol i.
\]

By Lemma~\ref{lem:whenequiv}, $p_1$ and $p_2$ are equivalent as elements of $\tilde S(D)$ if and only if:
\[
 p_1 - p_2 = \D^T(z_1,\dots, z_{n})^T
\]
for some $(z_1,\dots z_{n}) \in \Z^{n}$. 

Let $s_i$ be the restriction of the $i^{th}$ row of $\D$ to the first $r$ entries and let $\widehat s_i$ be the restriction of the $i^{th}$ row of $\D$ to the last $n-r$ entries. Then, the first $r$ entries of
\[
\D^T(z_1,\dots, z_{n})^T
\]
are given by
\[
 \sum_{i=1}^{n} z_is_i,
\]
and the last $n-r$ entries are given by
\[
\sum_{i=1}^{n} z_i\widehat s_i.
\]
From the structure of $\D$, $s_i$ and $\col i$ as well as $\widehat s_i$ and $\hcol i$ are closely related. In particular, for $i \in [1,r]$, we have $s_i = \col i$ and $\widehat s_i = -\hcol i$. For $i \in [r+1,n]$, we have $s_i = -\col i$ and $\widehat s_i = \hcol i$. 

This means that the first $r$ entries of
\[
\D^T(z_1,\dots z_r,-z_{r+1},\dots -z_{n})^T
\]
are given by
\[
\sum_{i=1}^{n} z_i \col i,
\]
and the last $(n-r)$ entries are given by
\[
\sum_{i=1}^{n} -z_i\hcol i.
\]
Hence the points $p_1$ and $p_2$ are equivalent as elements of $\tilde \Sand(D)$ if and only if we have:
\[
\sum_{i = 1}^{n}(d_i-z_i) \col i = \un 0 \hspace{ .4 cm} \text{ and } \hspace{ .4 cm} \sum_{i = 1}^{n}(\widehat d_i+z_i) \hcol i =\un 0.
\]
By the same logic that we used for Proposition~\ref{prop:nointersect}, the coefficients of the first sum form an element of $\ker_\R(D)$ while the coefficients of the second form an element of $\Ima_\R(D^T)$. Lemma~\ref{lem:orthogonality} again tells us that their dot product is 0. In other words:
\[
\sum_{i = 1}^{n}(d_i-z_i) \cdot (\widehat d_i + z_i) =0.
\]
For each $i$, there are 4 possibilities:

\begingroup
\setcounter{cdrthm}{0}
\renewcommand\thecdrthm{\arabic{cdrthm}}
\begin{enonce}[remark]{Case}\label{proof_def5.10_cas1} $i \in B_1 \cap B_2$:

$\widehat a_i = \widehat b_i = 0$, so $\widehat d_i = 0$. $0 < a_i,b_i <1$ so $-1 < d_i < 1$. If $z_i = 0$, then $(d_i-z_i) \cdot (\widehat d_i + z_i) = 0$. Otherwise, the two factors have a different sign and the product is negative. 
\end{enonce}

\begin{enonce}[remark]{Case}\label{proof_def5.10_cas2} $i \in B_1 \setminus B_2$:

$b_i = 0$, so $d_i = a_i$. This means that $0 < d_i <1$. $\widehat a_i = 0$, so $\widehat d_i = -\widehat b_i$. It follows that $-1 < \widehat d_i <0$. If $z_i > 0 $, then $d_i-z_i < 0$ and $\widehat d_i+ z_i > 0$. If $z_i \le 0$, then $d_i-z_i < 0$ and $\widehat d_i+ z_i > 0$. In either case, $(d_i-z_i) \cdot (\widehat d_i + z_i) < 0$.
\end{enonce}

\begin{enonce}[remark]{Case}\label{proof_def5.10_cas3} $i \in B_2 \setminus B_1$:

$a_i = 0$, so $d_i = -b_i$. This means that $-1 < d_i < 0$. $\widehat b_i = 0$, so $\widehat d_i = \widehat a_i$. It follows that $0 < \widehat d_i < 1$. If $z_i \ge 0 $, then $d_i-z_i < 0$ and $\widehat d_i+ z_i > 0$. If $z_i < 0$, then $d_i-z_i < 0$ and $\widehat d_i+ z_i > 0$. In either case, $(d_i-z_i) \cdot (\widehat d_i + z_i) < 0.$
\end{enonce}

\begin{enonce}[remark]{Case}\label{proof_def5.10_cas4} $i \not\in B_1 \cup B_2$:

$ a_i = b_i = 0$ so $d_i = 0$. $0 < \widehat a_i,\widehat b_i < 1$ so $-1 < \widehat d_i < 1$. If $z_i = 0$, then $(d_i-z_i) \cdot (\widehat d_i + z_i) = 0$. Otherwise, the two factors have a different sign and the product is negative. 
\end{enonce}
\setcounter{cdrthm}{10}
\endgroup


In all four cases the product is negative, unless we are always in case~\ref{proof_def5.10_cas1} or case~\ref{proof_def5.10_cas4} and $z_i = 0$ for all $i$. However, if $z_i = 0$ for all $i$, then $p_1 = p_2$. Thus, our claim holds by contradiction.

We showed that two distinct points $p_1$ and $p_2$ of $T(D)$ that are equivalent as elements of $\tilde \Sand(D)$ must each lie on the boundary of some $P(B)$. We now show by contradiction that they are both on the boundary of $T(D)$.

\looseness-1
Assume that $p_1$ is an interior point of $T(D)$. Since $T(D)$ is the union of non-degenerate parallelepipeds, there is some vector $w\in \R^n$ such that for all sufficiently small $\epsilon>0$, $p_2 + \epsilon w$ is in $T(D)$ but not on the boundary of any $P(B)$. If we make $\epsilon$ small enough, $p_1 + \epsilon w$ must be in $T(D)$ as well, since $p_1$ is an interior point of $T(D)$ by assumption. Moreover, $p_1 + \epsilon w$ and $p_2+ \epsilon w$ are equivalent as elements of $\tilde \Sand(D)$. We get a contradiction because both points are in $T(D)$, but $p_2+\epsilon w$ is not on the boundary of any $P(B)$. This means that $p_1$ and $p_2$ must both be on the boundary of $T(D)$. 
\end{proof}

The next corollary shows that copies of $T(D)$ can be used to periodically tile $\R^{n}$.

\begin{coro}\label{cor:tiling}
The set of translates $T(D) + \D^T(z_1,\dots, z_{n})^T$ for all $(z_1,\dots, z_{n}) \in \Z^{n}$ cover all of $\R^{n}$ and only intersect at their boundaries. 
\end{coro}

\begin{proof}
Consider any point $p\in \R^{n}$. By Lemma~\ref{lem:whenequiv}, the points which are equivalent to $p$ as elements of $\tilde S(D)$ are those of the form $p + \D^T(z_1,\dots, z_{n})$ for $(z_1,\dots, z_{n}) \in \Z^{n}$. Since these are exactly the translates of $T(D)$, the condition that the translates do not intersect except at their boundaries follows directly from Proposition~\ref{prop:lastDomino}.

We also have to show that the translates cover all of $\R^{n}$ given that they do not overlap except at their boundaries. We first note that $\Pi_\circ(\D^T)$ must tile $\R^{n}$ under the same translation because for every $p \in \R^{n}$, there is a unique solution to $(\D^T)p' = p$ (in particular $p' = (\D^T)^{-1} p$). We can map each point of $T(D)$ to a point in $\Pi_\circ(\D^T$) by translating it by an integer combination of columns of $\D^T$. Let $t$ be this piecewise translation from $T(D) \to \Pi_\circ(\D^T)$. Each translation preserves the volume of the region we transform and the only overlap is from the boundary of $T(D)$, which is a 0 volume set. It follows that the volume of the image of $t$ is equal to the volume of $T(D)$. Since $\Pi_\circ(\D^T$) has the same volume as $T(D)$, the set of points that are not in the image of $t$ must have volume $0$. 

Let $p$ be a point of $\Pi_\circ(\D^T$) that is not in the image of $t$. The preimage of $p$ is the collection of points in the same equivalence class with respect to $\tilde \Sand(D)$. By assumption, none of these points are in $T(D)$. Since $T(D)$ is closed, this means that none of these points are limit points of $T(D)$ either, so there is a neighborhood of $p$ that is also not in the image of $t$. However, this neighborhood must have positive volume, which is a contradiction. 
\end{proof}

\section{Constructing the Sandpile to Basis Multijections}\label{sec:multi}
In order to define our multijections, we will need $T(D)$ and an appropriate $\R^{n}$ direction vector.

\begin{defi}\label{def:shifting}
A \emph{shifting vector} $\w = (\w_1,\dots,\w_{n})$ of $D$ is a vector in $\R^{n}$ that is not in the span of a facet of $P(B)$ for any $B \in \B(D)$. 
\end{defi}


It will sometimes be useful to split our shifting vector into two smaller vectors. Consider the vectors $w = (w_1,\dots,w_r) \in \R^r$ and $\widehat w = (\widehat w_1,\dots, \widehat w_{n-r})\in \R^{n-r}$. We write $(w,\widehat w)$ for their concatenation, which is an $\R^n$ vector.
\begin{lemma}\label{lem:splitshift}
$(w,\widehat w)$ is a shifting vector for $D$ if and only if for all $B \in \B(D)$, $w$ does not lie in the span of any facet of $P_1(B)$ and $\widehat w$ does not lie in the span of any facet of $P_2(B)$. 
\end{lemma}
\begin{proof}
By definition, a point is in $P(B)$ if and only if it is in $P_1(B)$ when restricted to the first $r$ coordinates and $P_2(B)$ when restricted to the last $(n-r)$ coordinates. The lemma follows from the fact that $z+ \epsilon w$ is $(z_1,\dots, z_r) + \epsilon (w_1,\dots, w_r)$ when restricted to the first $r$ coordinates and $(\widehat z_1,\dots,\widehat z_{n-r}) + \epsilon (\widehat w_1,\dots, \widehat w_{n-r})$ when restricted to the last $(n-r)$ coordinates.
\end{proof}

\begin{defi}\label{def:reps}
Let $\w = (w,\widehat w)$ be a shifting vector.
\begin{itemize}
\item For any $v \in \R^{n}$, $ v$ is a \emph{$\w$-representative} of $\tilde \Sand(D)$ if $ v+\epsilon \w \in T(D)$ for all sufficiently small $\epsilon>0$. If $ v+ \epsilon \w \in P(B)$, we say that $ v$ is \emph{$\w$-associated} with $B$. 

\item For any $z \in \Z^{n}$, $z$ is a \emph{$\w$-representative} of $\Sand(D)$ if $z+\epsilon \w \in T(D)$ for all sufficiently small $\epsilon>0$. If $z+ \epsilon \w \in P(B)$, we say that $z$ is \emph{$\w$-associated} with $B$. 

\item For any $v \in \R^r$ if $v+ \epsilon w \in P_1(B)$ for all sufficiently small $\epsilon>0$, we say that $v$ is \emph{$w$-associated} with $B$. 

\item For any $\widehat v \in \R^{n-r}$ if $\widehat v+ \epsilon \widehat w \in P_2(B)$ for all sufficiently small $\epsilon>0$, we say that $\widehat v$ is \emph{$\widehat w$-associated} with $B$. 
\end{itemize}
\end{defi}

\begin{lemma}\label{lem:splitassoc}
Suppose $\w = (w,\widehat w)$ is a shifting vector, $v \in \R^r$, and $\widehat v \in \R^{n-r}$. Then, $(v,\widehat v)$ is $\w$-associated with $B$ if and only if $v$ is $w$-associated with $B$ and $\widehat v$ is $\widehat w$-associated with $B$. 
\end{lemma}
\begin{proof}
For any $\epsilon>0$, the first $r$ entries of $(v,\widehat v) + \epsilon \w$ are given by $v + \epsilon w$, and the last $n-r$ entries are given by $\widehat v + \epsilon \widehat w$. The lemma follows from the fact that $P_1(B)$ is $P(B)$ restricted to its first $r$ coordinates while $P_2(B)$ is $P(B)$ restricted to its last $n-r$ coordinates.
\end{proof}

\begin{lemma}\label{lem:associated}
Each $\w$-representative of $\tilde \Sand(D)$ or $\Sand(D)$ is $\w$-associated with exactly one $B \in \B(D)$. 
\end{lemma}
\begin{proof}
Since $\w$-representatives of $\Sand(D)$ are also $\w$-representatives of $\tilde \Sand(D)$, it suffices to prove the result for $\tilde \Sand(D)$. Let $p$ be a $\w$-representative of $\tilde \Sand(D)$. Because $T(D) = \bigcup P(B)$, we know that $p + \epsilon \w \in P(B)$ for some $B \in \B(D)$. Since $\w$ is not in the span of any facet of $P(B)$, $p + \epsilon \w$ must be in the interior of $P(B)$. By Proposition~\ref{prop:nointersect}, this is true for a unique $B$.
\end{proof}

\begin{prop}\label{prop:justone}
For any shifting vector $\w$, there is exactly one $\w$-representative in $\R^{n}$ for each equivalence class of $\tilde \Sand(D)$ and exactly one $\w$-representative in $\Z^{n}$ for each equivalence class of $\Sand(D)$.
\end{prop}

\begin{proof}
The second result is a direct corollary of the first (and could also be proven with an enumerative argument). By Corollary~\ref{cor:tiling}, every point $ p \in \R^{n}$ lies on some translation of $T(D)$ by an integer linear combination of the rows of $\D$. We can translate this point to a point on $T(D)$ without changing the equivalence class with respect to $\tilde \Sand(D)$. If $p$ maps to an interior point $p'$ of $T(D)$, then by Proposition~\ref{prop:lastDomino}, this is the unique point on $T(D)$ that is equivalent to $ p$. Furthermore, since $ p'$ is in the interior of $T(D)$, $ p'$ is always a $\w$-representative of $\tilde \Sand(D)$ regardless of $\w$.

If $ p$ maps to a boundary point of $T(D)$, then by Proposition~\ref{prop:lastDomino}, any point of $T(D)$ that is in the same $\tilde \Sand(D)$ equivalence class must also lie on the boundary of $T(D)$. Label these points as $\{ p_1,\dots, p_k\}$. We need to show that exactly one of these points is a $\w$-representative. 

By the condition that $\w$ is not in the span of any facet of $T(D)$, for all sufficiently small $\epsilon>0$, $ p_i + \epsilon \w$ must not lie on the boundary of $T(D)$ for any $i$. If $ p_i + \epsilon \w$ and $ p_j + \epsilon \w$ are both in $T(D)$ for $i \not= j$, then these are two distinct points in the interior of $T(D)$ that are equivalent as elements of $\tilde \Sand(D)$. This is impossible by Proposition~\ref{prop:lastDomino}. 

We have shown uniqueness, so we just need existence. Because $\w$ is not in the span of any facet of $T(D)$, we can choose $\epsilon>0$ so that all points between $p$ and $ p + \w\epsilon$ map to interior points of $T(D)$. Let $p'$ be the point mapped to by $ p + \w\epsilon$. Then, $p' - \epsilon \w$ must be equivalent to $p$ with respect to $\tilde \Sand(D)$. By our condition on $\epsilon$, we see that this point is a $\w$-representative. 
\end{proof}

\begin{prop}\label{prop:Ehrhart}
For any shifting vector $\w$, and for any $B \in \B(D)$, there are exactly $m(B)^2$ $\w$-representatives of $\Sand(D)$ that are $\w$-associated with $B$.
\end{prop}
To prove this result, we apply the following lemma from Ehrhart Theory:
\begin{lemma}[{\cite[Lemma~9.2]{Beck}}]\label{lem:Ehrhart}
For any integer matrix $M$, the number of integer points in the half-open fundamental parallelepiped $\Pi_\circ(M)$ is equal to its volume (the magnitude of $\det(M)$). 
\end{lemma}

\begin{proof}[Proof of Proposition~\ref{prop:Ehrhart}]
For some $B \in \B(D)$, let $\{\colB 1,\dots, \colB r\}$ be the columns of $D$ corresponding to $B$. Decompose $\w$ into the pair $(w,\widehat w)$ with $w \in \R^r$ and $\widehat w \in \R^{n-r}$. A point $v \in P_1(B)$ can be written as
\[
v = \sum_{i=1}^{r} a_i\colB i,
\]
with $0 \le a_i \le 1$ for all $i$. Because the $\colB i$ are linearly independent (otherwise $B$ would not be a basis), there is a unique way to write $w$ in the form:
\[
 w = \sum_{i=1}^{r} b_i\colB i,
\]
such that each $b_i \in \R$. By Lemma~\ref{lem:splitshift}, $w$ is not in the span of any facet of $P(B)$. This means that that $b_i \not= 0$ for all $i\in [r]$. For any $\epsilon \in \R$, we have:
\[
v+\epsilon w = \sum_i^{n} (a_i + \epsilon b_i)\colB i.
\]

From here, we see that $v$ is $w$-associated with $B$ if and only if $a_i \in (0,1]$ for $b_i < 0$ and $a_i \in [0,1)$ for $b_i > 0$. This region is the integer translation of a half-open fundamental parallelepiped with volume equal to the volume of $P_1(B)$. By an analogous line of reasoning, the points which are $\widehat w$-associated with $B$ form the integer translation of a half-open fundamental parallelepiped with volume equal to the volume of $P_2(B)$. It follows that the set of points that are $\w$-associated with $B$ is the direct product of these two regions: the integer translate of a half open parallelepiped with volume equal to the volume of $P(B)$.

By Lemma~\ref{lem:Ehrhart}, the number of integer points in this region is equal to this volume, and the integer translation does not change the number of integer points. Finally, by Lemma~\ref{lem:pipedvol}, the volume is $m(B)^2$, completing the proof. 
\end{proof}

We now define a function $\tilde f_{\w}$ from $\tilde \Sand(D) \to \B(D)$ given a shifting vector $\w$. For any $s \in \tilde \Sand(D)$, we first take the $\w$-representative $z$ of $s$ (which is unique by Proposition~\ref{prop:justone}). Then, we let $\tilde f_{\w}(s) = B$, where $B$ is the $\w$-associated basis of $z$ (which is unique by Lemma~\ref{lem:associated}). 

\begin{defi}\label{def:restrictedf}
$f_{\w}$ is $\tilde f_{\w}$ (as defined above) but with its domain restricted to $\Sand(D)$. 
\end{defi}

The following theorem is the main result of this paper.


\begin{theorem}\label{thm:Hellyeah}
For any $B\in \B(D)$, we have $|f_{\w}^{-1}(B)| = m(B)^2$.
\end{theorem}

\begin{proof}
We showed in Propositions~\ref{prop:justone} and~\ref{prop:nointersect} that $f_{\w}$ is a well-defined map from $\Sand(D)$ to $\B(D)$. The fact that $|f_{\w}^{-1}(B)| = m(B)^2$ is a corollary of Proposition~\ref{prop:Ehrhart}.
\end{proof}

\begin{thebibliography}{99}
\bibitem[4]{Beck} Matthias Beck and Sinai Robins, \emph{Computing the continuous discretely: Integer-point enumeration in polyhedra}, Undergraduate Texts in Mathematics, Springer, New York,2007.
\bibitem[11]{CutsFlows} Art M. Duval, Caroline J. Klivans and Jeremy L. Martin, ``Cuts and flows of cell complexes'', \emph{J. Algebraic Combin.}\textbf{41}(2015),no.4,969--999.
\bibitem[14]{Godsil} Chris Godsil and Gordon Royle, \emph{Algebraic graph theory}, Graduate Texts in Mathematics207, Springer-Verlag New York,2001.
\bibitem[20]{mythesis} Alex McDonough, \emph{Higher-dimensional sandpile groups and matrix-tree multijections}, Ph.D.thesis, Brown University,2021.
\end{thebibliography}
\end{document}
